\chapter{Stellar Temperature and Gas Statistics}

\section{Characteristic Stellar Temperatures}

\subsection{Internal Temperature Scale}
\begin{theory}
Using hydrostatic and virial scaling for an ideal-gas dominated star,
\begin{equation}
    k_B T \sim \frac{1}{3}\frac{GM\mu m_u}{R}.
\end{equation}
Typical stellar cores satisfy $T_I\sim 10^6$--$10^7\,$K.
\end{theory}

\subsection{Effective Surface Temperature}
\begin{theory}
From Stefan--Boltzmann law,
\begin{equation}
    L=4\pi R^2\sigma T_E^4,
    \qquad
    T_E=\left(\frac{L}{4\pi R^2\sigma}\right)^{1/4},
\end{equation}
with typical values $T_E\sim10^3$--$10^4\,$K.
\end{theory}

\begin{remark}
The large gap between $T_I$ and $T_E$ is a direct consequence of optical depth: radiation random-walks through dense stellar matter before escaping.
\end{remark}

\section{Validity of Ideal-Gas Approximation}

\subsection{Inter-particle Distance and Coulomb Scale}
\begin{foundation}
Average spacing is
\begin{equation}
    d \approx \left(\frac{\mu m_u}{\rho}\right)^{1/3}.
\end{equation}
A rough Coulomb scale is
\begin{equation}
    E_C \sim \frac{Z^2e^2}{4\pi\epsilon_0 d}.
\end{equation}
When thermal energy exceeds this scale, ideal-gas treatment is often sufficient.
\end{foundation}

\subsection{Virial Estimate for Thermal Energy}
\begin{theory}
For a self-gravitating system,
\begin{equation}
    2E_{\text{kin}}+E_{\text{grav}}=0,
\end{equation}
and for ideal gas,
\begin{equation}
    E_{\text{kin}}=\frac{3}{2}\frac{M}{\mu m_u}k_B T.
\end{equation}
Combined with $E_{\text{grav}}\sim -GM^2/R$ gives the standard central-temperature estimate.
\end{theory}

\section{Classical and Relativistic Ideal Gases}

\subsection{Energy--Pressure Relations}
\begin{theory}
General kinetic form:
\begin{equation}
    P=\frac{1}{3}n\langle\vec p\cdot\vec v\rangle.
\end{equation}
Hence
\begin{align}
    P &= \frac{2}{3}u \quad (\text{NR}), \\
    P &= \frac{1}{3}u \quad (\text{UR}).
\end{align}
\end{theory}

\subsection{Maxwell--Boltzmann Limit}
\begin{theory}
For classical non-relativistic particles,
\begin{equation}
    f(p)=A\exp\left(-\frac{p^2}{2mk_B T}\right),
\end{equation}
which leads to the ideal-gas law:
\begin{equation}
    PV=Nk_B T.
\end{equation}
The same equation of state also holds for ultra-relativistic classical gas, though with different energy scaling.
\end{theory}

\section{Quantum Statistics and Degeneracy}

\subsection{Fermi--Dirac and Bose--Einstein Distributions}
\begin{theory}
\begin{align}
    f_{\text{FD}}(\epsilon) &= \frac{1}{e^{(\epsilon-\mu)/k_BT}+1}, \\
    f_{\text{BE}}(\epsilon) &= \frac{1}{e^{(\epsilon-\mu)/k_BT}-1}.
\end{align}
Classical limit occurs when $e^{(\epsilon-\mu)/k_BT}\gg1$.
\end{theory}

\subsection{Quantum Criterion}
\begin{theory}
Define quantum concentration
\begin{equation}
    n_Q=g_s\left(\frac{2\pi mk_B T}{h^2}\right)^{3/2}.
\end{equation}
Then:
\begin{itemize}
    \item Classical regime: $n\ll n_Q$
    \item Quantum regime: $n\gtrsim n_Q$
\end{itemize}
\end{theory}

\begin{importantnote}
Degenerate pressure can dominate even when temperature is low.
It is set primarily by density through phase-space filling, not by thermal agitation.
\end{importantnote}

\subsection{Degenerate Fermi Gas at $T\to 0$}
\begin{theory}
At $T=0$, states are filled up to $p_F$:
\begin{equation}
    p_F=\hbar(3\pi^2 n)^{1/3} \quad (g_s=2).
\end{equation}
Degenerate pressure scales as
\begin{align}
    P_{\text{NR,deg}} &\propto n^{5/3}\propto \rho^{5/3}, \\
    P_{\text{UR,deg}} &\propto n^{4/3}\propto \rho^{4/3}.
\end{align}
\end{theory}

\begin{example}{Why white dwarfs shrink as mass increases}{ex:wd-mr}
In the non-relativistic degenerate regime, equilibrium between gravity and
$P\propto\rho^{5/3}$ leads to the approximate mass--radius trend
$R\propto M^{-1/3}$.

\smallskip\textbf{Interpretation.}
Adding mass increases gravity faster than degeneracy support can maintain the same size, so the star compresses.
\end{example}

\begin{extra}
The transition from $\rho^{5/3}$ to $\rho^{4/3}$ support explains why there is a maximum white-dwarf mass (Chandrasekhar limit).
\end{extra}
